\documentclass[../../../include/open-logic-section]{subfiles}

\begin{document}

\olfileid{sth}{replacement}{refproofs}
\olsection{ملحق: نتائج تتعلق بالاستبدال}

نثبت في هذا القسم الانعكاس داخل~$\ZF$، ثم نبين وجهًا يصير به الانعكاس
مكافئًا للاستبدال، ونختم بنتيجة لافتة في قوتهما. \emph{تنبيه: هذا
القسم، بفارق ظاهر، أرفع أجزاء الكتاب منزلة في الصعوبة الرياضية.}

ونقدم لمّة نستعمل فيها، اختصارًا، وسيلة ترميزية هي \emph{الخط
العلوي} للدلالة على متتاليات المتغيرات أو الأشياء. فالعبارة
«$\overline{\OLId{latin-ordinary}{a}}_\OLId{latin-ordinary}{k}$» تغني عن «$\OLId{latin-ordinary}{a}_{\OLId{latin-ordinary}{k}_\OLMathNumeral{1}}$، \dots، $\OLId{latin-ordinary}{a}_{\OLId{latin-ordinary}{k}_\OLId{latin-ordinary}{n}}$»، ويعين
السياق قيمة~$\OLId{latin-ordinary}{n}$.

\begin{lem}\ollabel{lemreflection}
لكل~$\OLMathNumeral{1} \leq \OLId{latin-ordinary}{i} \leq \OLId{latin-ordinary}{k}$، لتكن~$\phi_\OLId{latin-ordinary}{i}(\overline{\OLId{latin-ordinary}{v}}_\OLId{latin-ordinary}{i}, \OLId{latin-ordinary}{x})$ صيغة. عندئذ
يوجد لكل~$\OLId{greek-variable}{alpha}$ عدد~$\beta > \OLId{greek-variable}{alpha}$ بحيث، لأي
$\overline{\OLId{latin-ordinary}{a}}_\OLMathNumeral{1}, \ldots, \overline{\OLId{latin-ordinary}{a}}_\OLId{latin-ordinary}{k} \in \OLId{latin-looped}{V}_\beta$ ولكل
$\OLMathNumeral{1} \leq \OLId{latin-ordinary}{i} \leq \OLId{latin-ordinary}{k}$:
\[
	\exists \OLId{latin-ordinary}{x}\phi_\OLId{latin-ordinary}{i}(\overline{\OLId{latin-ordinary}{a}}_\OLId{latin-ordinary}{i}, \OLId{latin-ordinary}{x}) \rightarrow (\exists \OLId{latin-ordinary}{x} \in \OLId{latin-looped}{V}_\beta) \phi_\OLId{latin-ordinary}{i}(\overline{\OLId{latin-ordinary}{a}}_\OLId{latin-ordinary}{i}, \OLId{latin-ordinary}{x})
\]
\end{lem}

\begin{proof}
نعرّف حدًا~$\mu$ كما يلي: نجعل
$\mu(\overline{\OLId{latin-ordinary}{a}}_\OLMathNumeral{1}, \ldots, \overline{\OLId{latin-ordinary}{a}}_\OLId{latin-ordinary}{k})$ أول مرحلة~$\OLId{latin-looped}{V}$ تحقق،
لكل~$\OLMathNumeral{1} \leq \OLId{latin-ordinary}{i} \leq \OLId{latin-ordinary}{k}$، الشرطيات الآتية:
\[
\exists \OLId{latin-ordinary}{x}\phi_\OLId{latin-ordinary}{i}(\overline{\OLId{latin-ordinary}{a}}_\OLId{latin-ordinary}{i}, \OLId{latin-ordinary}{x}) \rightarrow (\exists \OLId{latin-ordinary}{x} \in \OLId{latin-looped}{V}) \phi_\OLId{latin-ordinary}{i}(\overline{\OLId{latin-ordinary}{a}}_\OLId{latin-ordinary}{i}, \OLId{latin-ordinary}{x})
\]
ولا يعسر التحقق من وجود~$\mu(\overline{\OLId{latin-ordinary}{a}}_\OLMathNumeral{1}, \ldots,
\overline{\OLId{latin-ordinary}{a}}_\OLId{latin-ordinary}{k})$ لكل~$\overline{\OLId{latin-ordinary}{a}}_\OLMathNumeral{1}, \ldots, \overline{\OLId{latin-ordinary}{a}}_\OLId{latin-ordinary}{k}$.
ثم نستخدم الاستبدال ومبرهنة العودية فنعرّف:
\begin{align*}
	\OLId{latin-looped}{S}_\OLMathNumeral{0} & = \OLId{latin-looped}{V}_{\OLId{greek-variable}{alpha}+\OLMathNumeral{1}}\\
	\OLId{latin-looped}{S}_{\OLId{latin-ordinary}{n}+\OLMathNumeral{1}} & = \OLId{latin-looped}{S}_\OLId{latin-ordinary}{n} \cup \bigcup
	\Setabs{\mu(\overline{\OLId{latin-ordinary}{a}}_\OLMathNumeral{1}, \ldots, \overline{\OLId{latin-ordinary}{a}}_\OLId{latin-ordinary}{k})}
	{\overline{\OLId{latin-ordinary}{a}}_\OLMathNumeral{1}, \ldots, \overline{\OLId{latin-ordinary}{a}}_\OLId{latin-ordinary}{k} \in \OLId{latin-looped}{S}_\OLId{latin-ordinary}{n}} \\
	\OLId{latin-looped}{S} &= \bigcup_{\OLId{latin-ordinary}{m} < \omega} \OLId{latin-looped}{S}_\OLId{latin-ordinary}{m}.
\end{align*}
كل~$\OLId{latin-looped}{S}_\OLId{latin-ordinary}{n}$، ومن ثم~$\OLId{latin-looped}{S}$ نفسها، مرحلة تلي~$\OLId{latin-looped}{V}_\OLId{greek-variable}{alpha}$. والآن ثبت
$\overline{\OLId{latin-ordinary}{a}}_\OLMathNumeral{1}$، \dots،~$\overline{\OLId{latin-ordinary}{a}}_\OLId{latin-ordinary}{k} \in \OLId{latin-looped}{S}$؛ فلا بد من
$\OLId{latin-ordinary}{n} < \omega$ بحيث~$\overline{\OLId{latin-ordinary}{a}}_\OLMathNumeral{1}$، \dots،
$\overline{\OLId{latin-ordinary}{a}}_\OLId{latin-ordinary}{k} \in \OLId{latin-looped}{S}_\OLId{latin-ordinary}{n}$. وثبت~$\OLMathNumeral{1} \leq \OLId{latin-ordinary}{i} \leq \OLId{latin-ordinary}{k}$، وافترض
$\exists \OLId{latin-ordinary}{x}\phi_\OLId{latin-ordinary}{i}(\overline{\OLId{latin-ordinary}{a}}_\OLId{latin-ordinary}{i},\OLId{latin-ordinary}{x})$. يقتضي البناء
$(\exists \OLId{latin-ordinary}{x} \in \mu(\overline{\OLId{latin-ordinary}{a}}_\OLMathNumeral{1},
\ldots, \overline{\OLId{latin-ordinary}{a}}_\OLId{latin-ordinary}{k}))\phi_\OLId{latin-ordinary}{i}(\overline{\OLId{latin-ordinary}{a}}_\OLId{latin-ordinary}{i}, \OLId{latin-ordinary}{x})$؛ ومنه
$(\exists \OLId{latin-ordinary}{x} \in \OLId{latin-looped}{S}_{\OLId{latin-ordinary}{n}+\OLMathNumeral{1}})\phi_\OLId{latin-ordinary}{i}(\overline{\OLId{latin-ordinary}{a}}_\OLId{latin-ordinary}{i}, \OLId{latin-ordinary}{x})$، ثم
$(\exists \OLId{latin-ordinary}{x} \in \OLId{latin-looped}{S})\phi_\OLId{latin-ordinary}{i}(\overline{\OLId{latin-ordinary}{a}}_\OLId{latin-ordinary}{i}, \OLId{latin-ordinary}{x})$. فهذه~$\OLId{latin-looped}{S}$ هي
$\OLId{latin-looped}{V}_\beta$ المطلوبة.
\end{proof}
\noindent
وبهذه اللمّة يصير إثبات
\olref[sth][replacement][ref]{reflectionschema} مباشرًا:

\begin{proof}[الإثبات]
ثبّت~$\OLId{greek-variable}{alpha}$. ولا يخل بالعموم أن نجعل روابط~$\phi$ مقصورة على
$\exists$ و$\lnot$ و$\land$، فهي وافية من جهة القوة التعبيرية.
ولتكن~$\psi_\OLMathNumeral{1}, \ldots, \psi_\OLId{latin-ordinary}{k}$ جميع الصيغ الجزئية لـ$\phi$، مرتبة
بحسب تعقيدها، بحيث~$\psi_\OLId{latin-ordinary}{k} = \phi$. ومن~\olref{lemreflection} يوجد
$\beta > \OLId{greek-variable}{alpha}$ بحيث، لأي~$\overline{\OLId{latin-ordinary}{a}}_\OLId{latin-ordinary}{i} \in \OLId{latin-looped}{V}_\beta$ ولكل
$\OLMathNumeral{1} \leq \OLId{latin-ordinary}{i} \leq \OLId{latin-ordinary}{k}$:
\begin{align}\label{reflectionnicelybehaved}
	\exists \OLId{latin-ordinary}{x}\psi_\OLId{latin-ordinary}{i}(\overline{\OLId{latin-ordinary}{a}}_\OLId{latin-ordinary}{i}, \OLId{latin-ordinary}{x}) \rightarrow
	(\exists \OLId{latin-ordinary}{x} \in \OLId{latin-looped}{V}_\beta) \psi_\OLId{latin-ordinary}{i}(\overline{\OLId{latin-ordinary}{a}}_\OLId{latin-ordinary}{i}, \OLId{latin-ordinary}{x})\tag{*}
\end{align}
ونثبت، بالاستقراء في تعقيد~$\psi_\OLId{latin-ordinary}{i}$، أنه لأي
$\overline{\OLId{latin-ordinary}{a}}_\OLId{latin-ordinary}{i} \in \OLId{latin-looped}{V}_\beta$:
$\psi_\OLId{latin-ordinary}{i}(\overline{\OLId{latin-ordinary}{a}}_\OLId{latin-ordinary}{i}) \leftrightarrow
\psi_\OLId{latin-ordinary}{i}^{\OLId{latin-looped}{V}_\beta}(\overline{\OLId{latin-ordinary}{a}}_\OLId{latin-ordinary}{i})$. أما إذا كانت~$\psi_\OLId{latin-ordinary}{i}$ ذرية
فالأمر ظاهر. ويثبت هذا التكافؤ كذلك أنه إذا كانت~$\psi_\OLId{latin-ordinary}{i}$ نفيًا أو
اقترانًا لصيغ جزئية تحققه، فإن~$\psi_\OLId{latin-ordinary}{i}$ نفسها تحققه؛ فلا تبقى إلا
حالة التكميم. ثبت
$\overline{\OLId{latin-ordinary}{a}}_\OLId{latin-ordinary}{i} \in \OLId{latin-looped}{V}_\beta$؛ فعندئذ:
\begin{align*}
	(\exists \OLId{latin-ordinary}{x} \psi_\OLId{latin-ordinary}{i}(\overline{\OLId{latin-ordinary}{a}}_\OLId{latin-ordinary}{i}, \OLId{latin-ordinary}{x}))^{\OLId{latin-looped}{V}_\beta}
	&\text{ إذا وفقط إذا }
	(\exists \OLId{latin-ordinary}{x} \in \OLId{latin-looped}{V}_\beta)\psi_\OLId{latin-ordinary}{i}^{\OLId{latin-looped}{V}_\beta}(\overline{\OLId{latin-ordinary}{a}}_\OLId{latin-ordinary}{i}, \OLId{latin-ordinary}{x})
	&&\text{بحكم التعريف}\\
	&\text{ إذا وفقط إذا }
	(\exists \OLId{latin-ordinary}{x} \in \OLId{latin-looped}{V}_\beta)\psi_\OLId{latin-ordinary}{i}(\overline{\OLId{latin-ordinary}{a}}_\OLId{latin-ordinary}{i},  \OLId{latin-ordinary}{x})
	&&\text{بحكم فرض الاستقراء}\\
	&\text{ إذا وفقط إذا }
	\exists \OLId{latin-ordinary}{x} \psi_\OLId{latin-ordinary}{i}(\overline{\OLId{latin-ordinary}{a}}_\OLId{latin-ordinary}{i}, \OLId{latin-ordinary}{x})
	&&\text{بموجب \eqref{reflectionnicelybehaved}}
\end{align*}
فتم الاستقراء، والنتيجة تابعة من~$\psi_\OLId{latin-ordinary}{k} = \phi$.
\end{proof}

ثبت إذن الانعكاس في~$\ZF$، وبرهاننا في جوهره طريق
\citet{Montague1961}. ونريد الآن أن نثبت، داخل~$\Z$، أن الانعكاس
يستلزم الاستبدال. والبرهان الآتي تبسيط لطريق~\citet{Levy1960}.

ولأننا نعمل في~$\Z$، لا يمكننا عرض الانعكاس بالصورة السابقة؛ فقد
استعملت تلك الصورة الرمز «$\OLId{latin-looped}{V}_\OLId{greek-variable}{alpha}$»، وهو غير قابل للتعريف في~$\Z$
(انظر~\olref[sth][spine][zf]{sec}). فنستعيض عنها بصورة تبدو أضعف:

\begin{defish}
\emph{الانعكاس الضعيف.} لكل صيغة~$\phi$، توجد مجموعة متعدية~$\OLId{latin-looped}{S}$ بحيث
تكون~$\OLMathNumeral{0}$ و$\OLMathNumeral{1}$ وأي وسائط لـ$\phi$ من !!{element}s~$\OLId{latin-looped}{S}$، ويكون
$(\forall \overline{\OLId{latin-ordinary}{x}} \in \OLId{latin-looped}{S})(\phi \liff \phi^\OLId{latin-looped}{S})$.
\end{defish}

ولكي نستعمل هذا في إثبات الاستبدال، نتبع أولًا
\citet[first part of Theorem 2]{Levy1960} ونثبت إمكان «عكس» صيغتين
معًا:

\begin{lem}[في $\Z + \text{الانعكاس الضعيف}$.]\ollabel{lem:reflect}
لكل صيغتين~$\psi, \chi$، توجد مجموعة متعدية~$\OLId{latin-looped}{S}$ بحيث يكون~$\OLMathNumeral{0}$ و$\OLMathNumeral{1}$
(وأي وسائط للصيغتين) من !!{element}s~$\OLId{latin-looped}{S}$، ويكون
$(\forall \overline{\OLId{latin-ordinary}{x}} \in \OLId{latin-looped}{S})((\psi \liff \psi^\OLId{latin-looped}{S}) \land (\chi
\liff \chi^\OLId{latin-looped}{S}))$.
\end{lem}

\begin{proof}
ضع~$\phi$ مساوية للصيغة~$(\OLId{latin-ordinary}{z} = \OLMathNumeral{0} \land \psi) \lor (\OLId{latin-ordinary}{z} = \OLMathNumeral{1} \land \chi)$.

وفي هذا اختصار ينبغي كشفه: فالعبارة «$\OLId{latin-ordinary}{z} = \OLMathNumeral{0}$» حقها أن تفصل إلى
«$\forall \OLId{latin-ordinary}{t}\, \OLId{latin-ordinary}{t} \notin \OLId{latin-ordinary}{z}$»، والعبارة «$\OLId{latin-ordinary}{z} =\OLMathNumeral{1}$» إلى «$\forall \OLId{latin-ordinary}{s}(\OLId{latin-ordinary}{s} \in \OLId{latin-ordinary}{z}
\liff \forall \OLId{latin-ordinary}{t}\, \OLId{latin-ordinary}{t} \notin \OLId{latin-ordinary}{s})$». غير أن~$\OLMathNumeral{0}, \OLMathNumeral{1} \in \OLId{latin-looped}{S}$، وأن~$\OLId{latin-looped}{S}$
متعدية؛ فهاتان الصيغتان \emph{مطلقتان} بالنسبة إلى~$\OLId{latin-looped}{S}$، أي تصدقان في
الشيء نفسه، سواء قيدنا مكمماتهما بـ$\OLId{latin-looped}{S}$ أم أطلقناها.
\footnote{وبصيغة أدق، إذا كانت~$\xi$ إحدى هاتين الصيغتين، فإن
$\xi(\OLId{latin-ordinary}{z}) \liff \xi^\OLId{latin-looped}{S}(\OLId{latin-ordinary}{z})$.}

ويعطينا الانعكاس الضعيف~$\OLId{latin-looped}{S}$ ملائمة بحيث:
\begin{align*}
	(\forall \OLId{latin-ordinary}{z}, \overline{\OLId{latin-ordinary}{x}} \in \OLId{latin-looped}{S})(&\phi \liff \phi^\OLId{latin-looped}{S})\\
	\text{أي }(\forall \OLId{latin-ordinary}{z}, \overline{\OLId{latin-ordinary}{x}} \in \OLId{latin-looped}{S})(&((\OLId{latin-ordinary}{z} = \OLMathNumeral{0} \land \psi) \lor (\OLId{latin-ordinary}{z} = \OLMathNumeral{1} \land \chi)) \liff {}\\
	&\phantom{(}((\OLId{latin-ordinary}{z} = \OLMathNumeral{0} \land \psi) \lor (\OLId{latin-ordinary}{z} = \OLMathNumeral{1} \land \chi))^\OLId{latin-looped}{S})\\
	\text{أي }(\forall \OLId{latin-ordinary}{z}, \overline{\OLId{latin-ordinary}{x}} \in \OLId{latin-looped}{S})(&((\OLId{latin-ordinary}{z} = \OLMathNumeral{0} \land \psi) \lor (\OLId{latin-ordinary}{z} = \OLMathNumeral{1} \land \chi))\liff {}\\
	&\phantom{(}((\OLId{latin-ordinary}{z} = \OLMathNumeral{0} \land \psi^\OLId{latin-looped}{S}) \lor (\OLId{latin-ordinary}{z} = \OLMathNumeral{1} \land \chi^\OLId{latin-looped}{S})))\\
	\text{أي }(\forall \overline{\OLId{latin-ordinary}{x}} \in \OLId{latin-looped}{S})(&(\psi \liff \psi^\OLId{latin-looped}{S}) \land (\chi \liff \chi^\OLId{latin-looped}{S}))
\end{align*}
فالادعاء الثاني يقتضي الثالث لإطلاق «$\OLId{latin-ordinary}{z} = \OLMathNumeral{0}$» و«$\OLId{latin-ordinary}{z}=\OLMathNumeral{1}$» بالنسبة
إلى~$\OLId{latin-looped}{S}$؛ والرابع تابع لأن~$\OLMathNumeral{0} \neq \OLMathNumeral{1}$.
\end{proof}\noindent
ثم نحصل على الاستبدال باتباع
\citet[Theorem 6]{Levy1960} مع تبسيط الحجة:

\begin{thm}[في $\Z$ + الانعكاس الضعيف]\label{thm:replacement}
لكل صيغة~$\phi(\OLId{latin-ordinary}{v},\OLId{latin-ordinary}{w})$ ولكل~$\OLId{latin-looped}{A}$، إذا كان
$(\forall \OLId{latin-ordinary}{x} \in \OLId{latin-looped}{A})\lexists![\OLId{latin-ordinary}{y}][\phi(\OLId{latin-ordinary}{x},\OLId{latin-ordinary}{y})]$، فإن
$\Setabs{\OLId{latin-ordinary}{y}}{(\exists \OLId{latin-ordinary}{x} \in \OLId{latin-looped}{A})\phi(\OLId{latin-ordinary}{x},\OLId{latin-ordinary}{y})}$ موجودة.
\end{thm}

\begin{proof}
ثبّت~$\OLId{latin-looped}{A}$ بحيث~$(\forall \OLId{latin-ordinary}{x} \in \OLId{latin-looped}{A})\lexists![\OLId{latin-ordinary}{y}][\phi(\OLId{latin-ordinary}{x},\OLId{latin-ordinary}{y})]$، وعرّف
الصيغتين:
\begin{align*}
	\psi &\text{ هي } (\phi(\OLId{latin-ordinary}{x}, \OLId{latin-ordinary}{z}) \land \OLId{latin-looped}{A} = \OLId{latin-looped}{A})\\
	\chi &\text{ هي } \lexists[\OLId{latin-ordinary}{y}][\phi(\OLId{latin-ordinary}{x}, \OLId{latin-ordinary}{y})]
\end{align*}
من~\olref{lem:reflect}، ولأن~$\OLId{latin-looped}{A}$ وسيط في~$\psi$، توجد مجموعة متعدية
$\OLId{latin-looped}{S}$ بحيث~$\OLMathNumeral{0}, \OLMathNumeral{1}, \OLId{latin-looped}{A} \in \OLId{latin-looped}{S}$، ومعها سائر الوسائط، وبحيث:
\[
	(\forall \OLId{latin-ordinary}{x},\OLId{latin-ordinary}{z} \in \OLId{latin-looped}{S})((\psi \liff \psi^\OLId{latin-looped}{S}) \land (\chi \liff \chi^\OLId{latin-looped}{S}))
\]
فيكون، على الخصوص:
\begin{align*}
	(\forall  \OLId{latin-ordinary}{x}, \OLId{latin-ordinary}{z} \in \OLId{latin-looped}{S})(&\phi(\OLId{latin-ordinary}{x}, \OLId{latin-ordinary}{z}) \liff \phi^\OLId{latin-looped}{S}(\OLId{latin-ordinary}{x}, \OLId{latin-ordinary}{z}))\\
	(\forall \OLId{latin-ordinary}{x} \in \OLId{latin-looped}{S})(&\exists \OLId{latin-ordinary}{y}\phi(\OLId{latin-ordinary}{x}, \OLId{latin-ordinary}{y}) \liff (\exists \OLId{latin-ordinary}{y} \in \OLId{latin-looped}{S})\phi^\OLId{latin-looped}{S}(\OLId{latin-ordinary}{x}, \OLId{latin-ordinary}{y}))
\end{align*}
نجمع بينهما، ونلاحظ أن~$\OLId{latin-looped}{A} \subseteq \OLId{latin-looped}{S}$ لأن~$\OLId{latin-looped}{A} \in \OLId{latin-looped}{S}$ و$\OLId{latin-looped}{S}$ متعدية،
فنحصل على:
\begin{align*}
	(\forall \OLId{latin-ordinary}{x} \in \OLId{latin-looped}{A})(&\exists \OLId{latin-ordinary}{y}\phi(\OLId{latin-ordinary}{x}, \OLId{latin-ordinary}{y}) \liff (\exists \OLId{latin-ordinary}{y} \in \OLId{latin-looped}{S})\phi(\OLId{latin-ordinary}{x}, \OLId{latin-ordinary}{y}))
\end{align*}
ولما كان~$(\forall \OLId{latin-ordinary}{x} \in \OLId{latin-looped}{A})\lexists![\OLId{latin-ordinary}{y} \in \OLId{latin-looped}{S}][\phi(\OLId{latin-ordinary}{x},\OLId{latin-ordinary}{y})]$، من
الفرض~$(\forall \OLId{latin-ordinary}{x} \in \OLId{latin-looped}{A})\lexists![\OLId{latin-ordinary}{y}][\phi(\OLId{latin-ordinary}{x}, \OLId{latin-ordinary}{y})]$، أعطانا الفصل
$\Setabs{\OLId{latin-ordinary}{y} \in \OLId{latin-looped}{S}}{(\exists \OLId{latin-ordinary}{x} \in \OLId{latin-looped}{A}) \phi(\OLId{latin-ordinary}{x}, \OLId{latin-ordinary}{y})} = \Setabs{\OLId{latin-ordinary}{y}}{(\exists
\OLId{latin-ordinary}{x} \in \OLId{latin-looped}{A}) \phi(\OLId{latin-ordinary}{x}, \OLId{latin-ordinary}{y})}$.
\end{proof}

\end{document}
